\chapter*{Kurzfassung}
\addcontentsline{toc}{chapter}{Kurzfassung}

In schulischen und beruflichen Teamprojekten sind Aufgabenverwaltung und Kommunikation
häufig auf mehrere getrennte Werkzeuge verteilt, was zu Informationsverlust und
sogenanntem \textit{Context Switching} führt. Die vorliegende Diplomarbeit beschreibt
Entwicklung und Implementierung von \textit{EffSync}, einer webbasierten Anwendung, die
Aufgabenverwaltung und Echtzeitkommunikation in einem gemeinsamen System zusammenführt.

Das System wurde mit \textit{Vue.js} im Frontend sowie \textit{NestJS}, \textit{Prisma}
und \textit{PostgreSQL} im Backend umgesetzt. Im Mittelpunkt stehen eine hierarchische Task-Struktur mit serverseitig abgesicherter Erkennung wechselseitiger Abhängigkeiten, ein kontextbasiertes Berechtigungskonzept sowie ein nativer, WebSocket-basierter Live-Chat. Die Qualität der
Umsetzung wurde durch eine mehrstufige Teststrategie (Unit-, API- und End-to-End-Tests)
sowie einen moderierten Usability-Test abgesichert.

Die Ergebnisse belegen, dass sämtliche für das Semester festgelegten SMART-Ziele (spezifisch, messbar, attraktiv, realistisch und terminiert) vollständig erreicht oder sogar in Teilen übertroffen wurden. Den Abschluss der Arbeit bildet eine kritische Reflexion der gewählten Backend-Architektur sowie ein Ausblick auf mögliche Erweiterungen wie Mandantenfähigkeit, horizontale Skalierung und native Mobile-Apps.
\vspace{1cm}

\chapter*{Abstract}
\addcontentsline{toc}{chapter}{Abstract}

In academic and professional team projects, task management and communication are
frequently distributed across multiple separate tools, leading to information loss and
so-called \textit{context switching}. This diploma thesis describes the development and
implementation of \textit{EffSync}, a web-based application that unifies task management
and real-time communication within a single system. The system was implemented using
\textit{Vue.js} on the frontend and \textit{NestJS}, \textit{Prisma}, and \textit{PostgreSQL}
on the backend. At its core are a hierarchical task structure with server-side validated
detection of mutual dependencies, a context-based permission concept, and a native,
WebSocket-based live chat. The quality of the implementation was ensured through a
multi-stage testing strategy (unit, API, and end-to-end tests) as well as a moderated
usability test.
The results demonstrate that all SMART goals (specific, measurable, attractive, realistic
and time-bound) defined for the semester were fully achieved or in some cases even
exceeded. The thesis concludes with a critical reflection on the chosen backend
architecture and an outlook on possible extensions such as multi-tenancy, horizontal
scaling, and native mobile apps.
